% Generated draft: requires mathematical and source review.
\documentclass[11pt]{article}
\usepackage[margin=25mm]{geometry}
\usepackage{fontspec}
\setmainfont{Latin Modern Roman}
\usepackage{amsmath,amssymb,mathtools,unicode-math}
\setmathfont{Latin Modern Math}
\usepackage{enumitem,needspace,xurl,hyperref}
\setlength{\parindent}{0pt}
\setlength{\parskip}{5pt}
\setlength{\emergencystretch}{3em}
\newcommand{\N}{\mathbb{N}}
\newcommand{\Z}{\mathbb{Z}}
\newcommand{\Q}{\mathbb{Q}}
\newcommand{\R}{\mathbb{R}}
\newcommand{\C}{\mathbb{C}}
\newcommand{\F}{\mathbb{F}}
\newcommand{\A}{\mathbb{A}}
\newcommand{\PP}{\mathbb{P}}
\newcommand{\E}{\mathbb{E}}
\newcommand{\eps}{\varepsilon}
\newcommand{\dd}{\,\mathrm{d}}
\begin{document}
\section*{UnsolvedMath ID: 20000563}
\textbf{Source problem code:} catalogue entry
\subsection{Definitions and mathematical statement}
The source asks how many ``good partitions'' the proof produces and calls for more quantitative versions. It does not name the proof or define goodness in this item. The attached aid interprets this as counting Weaver-good partitions for the Marcus--Spielman--Srivastava theorem, which supplies the following precise possible target.

Let \(v_1,\ldots,v_m\in\C^d\), \(d\geq1\), be a rank-one Parseval frame:
\[
A_i=v_iv_i^*,\qquad \sum_{i=1}^m A_i=I_d,\qquad
\|v_i\|^2\leq\alpha.
\]
For \(S\subset\{1,\ldots,m\}\), use signed discrepancy
\[
D(S)=\left\|\sum_{i\in S}A_i-\sum_{i\notin S}A_i\right\|_{\mathrm{op}}
=\left\|2\sum_{i\in S}A_i-I_d\right\|_{\mathrm{op}}.
\]
For a threshold \(\delta\geq0\), call the unordered bipartition \(\{S,S^c\}\) good if \(D(S)\leq\delta\), and let \(N_{\mathrm{good}}\) be their number. Empty parts are allowed in this explicit convention. The counting question asks for lower bounds for \(N_{\mathrm{good}}\), or for the fraction \(N_{\mathrm{good}}/2^{m-1}\), in terms of \(d,m,\alpha,\delta\), uniformly over the permitted frames. Labeled subsets instead give twice this count. The normalization and labeled/unordered choice are not supplied in the original item.

The attached aid reports the specific theorem
\[
0<\alpha\leq\tfrac14,\qquad
\delta=\sqrt{8\alpha}+2\alpha
\quad\Longrightarrow\quad N_{\mathrm{good}}\geq m+1.
\]
The discrepancy normalization is the signed one above; an estimate for deviation of \(\sum_{i\in S}A_i\) from \(I_d/2\) has half that threshold. The aid calls this the historical MSS bound. It explains its claimed abundance through a sharper Bownik--Casazza--Marcus--Speegle witness with strict one-flip slack and a weighted Hamming-stability principle: changing the membership of indices in \(J\) changes signed discrepancy by at most \(2\sum_{i\in J}\|v_i\|^2\). The precise improved witness bound is not supplied, so it is not reconstructed as a new theorem here.

Other supplied partial results are a matrix-Bernstein counting bound, an exact repeated-coordinate product count, and an abstract example in which common interlacing alone certifies only one favorable leaf. Their full formulas and hypotheses are absent. The aid says the repeated-coordinate family has good fraction tending to zero at the certified Kadison--Singer thresholds, so no universal positive-fraction strengthening at that threshold follows. This does not contradict a worst-case lower count of \(m+1\), which can be a vanishing fraction of \(2^{m-1}\). Nor does the abstract interlacing example establish a universal obstruction to every more quantitative method. These are supplied unreviewed partial claims; they do not specify every possible interpretation of the source's request.
\subsection{Short English statement}
Quantify the number of partitions satisfying a Weaver-type discrepancy guarantee. The supplied aid gives an \(m+1\) lower bound for unordered bipartitions at one historical threshold, while reporting examples whose good fraction vanishes. Counting partitions and proving a positive proportion are different goals.
\subsection{Sources}
\begin{itemize}
\item AIM workshop \emph{Beyond Kadison--Singer: paving and consequences}, item b22: \url{https://docs.google.com/document/d/1XRw0WyG5PshqboCp99tovCRD1GCK0jXZiNLH_nw618c/edit}.
\item Local registry \texttt{data/unsolvedmath/problems.json}, permanent ID \(20000563\), exact source code \texttt{catalogue entry}; canonical \texttt{aim-analysis-notes.json}, notes[94].
\item The attached aid names Marcus--Spielman--Srivastava and Bownik--Casazza--Marcus--Speegle, and invokes matrix Bernstein and common interlacing, without full bibliographic references. Attempt 1 is explicitly machine-generated, not peer reviewed, and requires independent verification before citation.
\end{itemize}
\end{document}
